\documentclass[twocolumn,10pt,a4paper]{article}

% Page layout
\usepackage[margin=2cm, columnsep=0.6cm]{geometry}

% Essential packages
\usepackage{amsmath,amssymb,amsthm}
\usepackage{mathtools}
\usepackage{bm}
\usepackage{graphicx}
\usepackage{hyperref}
\usepackage{physics}
\usepackage{booktabs}
\usepackage{array}
\usepackage{multirow}
\usepackage{enumitem}

% Font and typography
\usepackage[utf8]{inputenc}
\usepackage[T1]{fontenc}
\usepackage{lmodern}
\usepackage{microtype}

\usepackage{caption}
\captionsetup{
  font=small,          % or footnotesize, scriptsize
  labelfont=footnotesize,        % Keep "Figure 1" bold
  textfont=it          % Optional: italicize caption text
}

% Theorem environments
\newtheorem{theorem}{Theorem}[section]
\newtheorem{lemma}[theorem]{Lemma}
\newtheorem{proposition}[theorem]{Proposition}
\newtheorem{corollary}[theorem]{Corollary}
\theoremstyle{definition}
\newtheorem{definition}[theorem]{Definition}
\newtheorem{axiom}{Axiom}
\newtheorem{example}[theorem]{Example}
\theoremstyle{remark}
\newtheorem{remark}[theorem]{Remark}

% Hyperref setup
\hypersetup{
    colorlinks=true,
    linkcolor=blue,
    citecolor=blue,
    urlcolor=blue,
    pdftitle={On the Geometric Partitioning Consequences of Composition-Inflation in Bounded Phase Space: Exponential Categorical Resolution from Integer Compositions in d-Dimensional Entropy Space},
    pdfauthor={Kundai Farai Sachikonye}
}

% Custom commands
\newcommand{\kB}{k_{\mathrm{B}}}
\newcommand{\Hilbert}{\mathcal{H}}
\newcommand{\Partition}{\mathcal{P}}
\newcommand{\Cat}{\mathcal{C}}
\newcommand{\Ocat}{\hat{O}_{\mathrm{cat}}}
\newcommand{\Ophys}{\hat{O}_{\mathrm{phys}}}
\newcommand{\Scat}{S_{\mathrm{cat}}}
\newcommand{\Sk}{S_k}
\newcommand{\St}{S_t}
\newcommand{\Se}{S_e}
\newcommand{\tP}{t_{\mathrm{P}}}
\newcommand{\lP}{\ell_{\mathrm{P}}}
\newcommand{\nP}{n_{\mathrm{P}}}

\begin{document}

% Title
\title{\textbf{On the Geometric Partitioning Consequences of Composition-Inflation in Bounded Phase Space: Exponential Categorical Resolution from Integer Compositions in $d$-Dimensional Entropy Space}}

\author{
Kundai Farai Sachikonye\\
AIMe Registry for Artificial Intelligence\\
Experimental Bioinformatics\\
Maximus-von-Imhof-Forum 3\\
85354 Freising, Germany\\
\texttt{kundai.sachikonye@bitspark.com}
}

\date{\today}

\maketitle

%==============================================================================
% ABSTRACT
%==============================================================================

\begin{abstract}
\noindent We derive the \emph{composition-inflation mechanism}: for a bounded oscillatory system in $d$-dimensional entropy space with $n$ oscillation cycles, the number of categorically distinguishable trajectories is $T(n,d) = d \cdot (d+1)^{n-1}$, not~$n$. This exponential growth arises because (i) the $n$ cycles admit $2^{n-1}$ distinct compositions (ordered partitions), and (ii) each step in a composition is labeled by one of~$d$ entropy dimensions. For the three-dimensional S-entropy space $(\Sk, \St, \Se)$, $T(n,3) = 3 \cdot 4^{n-1}$. We prove that angular resolution $\Delta\theta = 2\pi / T(n,d)$ is dimensionless and therefore not subject to the Planck time constraint, which applies only to temporal intervals measured in seconds. We demonstrate that 56 oscillation cycles of caesium-133 suffice to generate enough distinguishable trajectories ($3.8 \times 10^{33}$) to match Planck-time angular resolution---not by ticking faster, but by exploiting the combinatorial richness of trajectory structure within the same ticks. This mechanism provides a unified derivation of trans-Planckian categorical resolution, subsuming the five enhancement mechanisms previously identified~\cite{sachikonye2025trans}, and establishes that the Planck time is not a fundamental limit on measurement but an artifact of the SI unit system's conflation of angular resolution with temporal duration.

\medskip
\noindent\textbf{Keywords:} composition-inflation, bounded phase space, integer compositions, ternary S-entropy encoding, categorical trajectory count, angular resolution, trans-Planckian measurement, Planck depth, dimensional labeling, Poincaré recurrence, partition propagation bound, categorical angular units, gravitational constant irreducibility, oscillation counting

\end{abstract}


%==============================================================================
\section{Introduction}
\label{sec:introduction}
%==============================================================================

\subsection{The Trans-Planckian Question}

The Planck time $\tP = \sqrt{\hbar G / c^5} \approx 5.391 \times 10^{-44}$~s is conventionally regarded as the shortest physically meaningful time interval~\cite{planck1901,garay1995}. The standard argument combines the Heisenberg uncertainty principle $\Delta E \cdot \Delta t \geq \hbar/2$ with general relativity: resolving a time interval $\delta t$ requires energy uncertainty $\Delta E \gtrsim \hbar / \delta t$, and for $\delta t \sim \tP$ this energy collapses spacetime into a black hole of Schwarzschild radius $r_S \geq 2\lP$. The measurement destroys itself.

This argument is correct for \emph{temporal intervals measured in seconds}. But it is silent on \emph{categorical trajectory counts}, which are dimensionless. The question ``can measurement resolve structures finer than $5.39 \times 10^{-44}$~s?'' is ill-posed as stated. Resolution measured in seconds inherits the Planck limit. Resolution measured in the number of distinguishable categorical trajectories does not.

\subsection{The Key Mechanism}

Consider a bounded oscillatory system---a pendulum, an atomic transition, a molecular vibration---that completes $n$ full cycles. The naive count of distinguishable temporal states is $n$: one state per cycle. But this dramatically undercounts the actual categorical structure.

A system with $n$ cycles can distribute those cycles among different phases according to $2^{n-1}$ \emph{compositions} (ordered partitions). Each step in a composition can be labeled by one of $d$ entropy dimensions. The total number of categorically distinguishable trajectories is therefore not $n$ but
\begin{equation}
\label{eq:main}
T(n,d) = d \cdot (d+1)^{n-1},
\end{equation}
which grows \emph{exponentially} in $n$.

For the three-dimensional S-entropy space with coordinates $(\Sk, \St, \Se)$, this yields $T(n,3) = 3 \cdot 4^{n-1}$. After merely 56 oscillation cycles of caesium-133, the trajectory count exceeds $10^{33}$---the number of Planck-time intervals in one caesium period.

\subsection{Outline}

This paper derives the composition-inflation formula from first principles (Sections~\ref{sec:compositions}--\ref{sec:labeling}), proves trajectory distinguishability (Section~\ref{sec:distinguishability}), establishes the angular resolution result (Section~\ref{sec:angular}), computes the Planck depth for physical oscillators (Section~\ref{sec:planck_depth}), demonstrates unification of the five previously identified enhancement mechanisms (Section~\ref{sec:unification}), develops the radian reformulation (Section~\ref{sec:radian}), explores consequences (Section~\ref{sec:consequences}), and discusses implications (Section~\ref{sec:discussion}).


%==============================================================================
\section{Integer Compositions}
\label{sec:compositions}
%==============================================================================

We begin with the combinatorial foundation. The theory of integer compositions is classical~\cite{andrews1976,heubach2009}, but its physical interpretation in the context of bounded oscillatory systems is new.

\begin{definition}[Composition]
\label{def:composition}
A \emph{composition} of a positive integer $n$ is an ordered sequence of positive integers $(c_1, c_2, \ldots, c_k)$ such that
\begin{equation}
c_1 + c_2 + \cdots + c_k = n, \quad c_i \geq 1 \;\;\forall\, i.
\end{equation}
The integer $k$ is called the \emph{number of parts} of the composition. Two compositions that differ only in the ordering of their parts are considered distinct.
\end{definition}

\begin{example}
The compositions of $n = 4$ are:
\begin{align*}
k=1: &\quad (4) \\
k=2: &\quad (3,1),\; (1,3),\; (2,2) \\
k=3: &\quad (2,1,1),\; (1,2,1),\; (1,1,2) \\
k=4: &\quad (1,1,1,1)
\end{align*}
yielding $1 + 3 + 3 + 1 = 8 = 2^3 = 2^{4-1}$ compositions in total.
\end{example}

\begin{theorem}[Composition Count]
\label{thm:comp_count}
The total number of compositions of $n$ is $2^{n-1}$.
\end{theorem}

\begin{proof}
Represent $n$ as a row of $n$ identical objects (stars). Between consecutive objects there are $n-1$ gaps. Each composition corresponds to a choice of which gaps contain dividers: a divider in gap $i$ separates the $i$-th and $(i+1)$-th objects into different parts. Each gap either contains a divider or does not, giving two independent choices per gap. The total number of configurations is $2^{n-1}$.

Formally, define the map $\varphi: \{0,1\}^{n-1} \to \mathcal{C}(n)$ from binary strings of length $n-1$ to compositions of $n$ by: given $b = (b_1, \ldots, b_{n-1}) \in \{0,1\}^{n-1}$, let $0 = i_0 < i_1 < \cdots < i_{k-1} < i_k = n$ be the positions where $b_{i_j} = 1$ (with $i_0 = 0$ and $i_k = n$ as sentinels), and set $c_j = i_j - i_{j-1}$. This map is a bijection: every composition determines a unique binary string, and conversely. Since $|\{0,1\}^{n-1}| = 2^{n-1}$, the result follows.
\end{proof}

\begin{theorem}[Compositions by Part Count]
\label{thm:comp_parts}
The number of compositions of $n$ into exactly $k$ parts is $\binom{n-1}{k-1}$.
\end{theorem}

\begin{proof}
A composition of $n$ into $k$ parts corresponds to placing exactly $k-1$ dividers among the $n-1$ gaps between $n$ objects. The number of ways to choose $k-1$ gaps from $n-1$ is $\binom{n-1}{k-1}$.

Alternatively, by the bijection $\varphi$ from Theorem~\ref{thm:comp_count}, a composition into $k$ parts corresponds to a binary string of length $n-1$ with exactly $k-1$ ones. The count is $\binom{n-1}{k-1}$.
\end{proof}

\begin{corollary}
\label{cor:sum_binom}
$\sum_{k=1}^{n} \binom{n-1}{k-1} = 2^{n-1}$.
\end{corollary}

\begin{proof}
Substituting $j = k-1$: $\sum_{j=0}^{n-1} \binom{n-1}{j} = 2^{n-1}$ by the binomial theorem applied to $(1+1)^{n-1}$.
\end{proof}

\begin{example}
\label{ex:n3}
For $n = 3$, the compositions are:
\begin{itemize}
\item $k=1$: $(3)$ --- $\binom{2}{0} = 1$ composition.
\item $k=2$: $(2,1),\; (1,2)$ --- $\binom{2}{1} = 2$ compositions.
\item $k=3$: $(1,1,1)$ --- $\binom{2}{2} = 1$ composition.
\end{itemize}
Total: $1 + 2 + 1 = 4 = 2^2 = 2^{3-1}$. \checkmark
\end{example}

\subsection{Physical Interpretation}

For a bounded oscillatory system with period $n$ ticks, each composition represents a different \emph{temporal structure}---a different way of dwelling in and transitioning between categorical states within the same total period.

The composition $(2,1)$ means: dwell for 2 ticks in one categorical state, then transition for 1 tick. The composition $(1,2)$ means: transition for 1 tick, then dwell for 2 ticks. These produce different trajectories through entropy space even though they cover the same total period of 3 ticks.

The key insight is that compositions are \emph{ordered}: $(2,1) \neq (1,2)$. The temporal sequence matters. A system that dwells first and transitions second is categorically distinct from one that transitions first and dwells second. This is the fundamental difference between compositions (where order matters) and partitions (where it does not). Bounded oscillatory systems preserve ordering information through their trajectory structure, making compositions the correct combinatorial object.


%==============================================================================
\section{Dimensional Labeling}
\label{sec:labeling}
%==============================================================================

The S-entropy space is a $d$-dimensional unit hypercube $\mathbf{S} = [0,1]^d$ with coordinates $(S_1, S_2, \ldots, S_d)$. For the physical case relevant to bounded molecular dynamics, $d = 3$ with coordinates $(\Sk, \St, \Se)$ representing kinetic, thermal, and electronic entropy dimensions respectively~\cite{sachikonye2025bounded}.

\begin{definition}[Labeled Composition]
\label{def:labeled_comp}
A \emph{labeled composition} of $n$ in $d$ dimensions is a pair $(\mathbf{C}, \mathbf{L})$ where:
\begin{itemize}
\item $\mathbf{C} = (c_1, c_2, \ldots, c_k)$ is a composition of $n$ into $k$ parts,
\item $\mathbf{L} = (\ell_1, \ell_2, \ldots, \ell_k)$ is a \emph{labeling} with $\ell_i \in \{1, 2, \ldots, d\}$,
\end{itemize}
where the label $\ell_i$ specifies which entropy dimension is being refined during the $i$-th part of the trajectory.
\end{definition}

The label carries physical meaning: $\ell_i = 1$ means the system is refining its position along $S_1$ (e.g., $\Sk$) during the $i$-th segment; $\ell_i = 2$ means refinement along $S_2$ (e.g., $\St$); and so forth. Different labelings of the same composition produce genuinely different trajectories through entropy space.

\begin{example}
For $n = 3$, $d = 3$, and the composition $\mathbf{C} = (2,1)$, the possible labelings are:
\begin{align*}
&(\Sk, \Sk),\; (\Sk, \St),\; (\Sk, \Se), \\
&(\St, \Sk),\; (\St, \St),\; (\St, \Se), \\
&(\Se, \Sk),\; (\Se, \St),\; (\Se, \Se).
\end{align*}
Each of these 9 labeled compositions describes a different trajectory: ``dwell 2 ticks in kinetic, then 1 tick in thermal'' differs from ``dwell 2 ticks in kinetic, then 1 tick in electronic.''
\end{example}

\begin{theorem}[Labeled Composition Count by Parts]
\label{thm:labeled_parts}
The number of labeled compositions of $n$ into exactly $k$ parts with $d$ dimension labels is
\begin{equation}
\binom{n-1}{k-1} \cdot d^k.
\end{equation}
\end{theorem}

\begin{proof}
By Theorem~\ref{thm:comp_parts}, there are $\binom{n-1}{k-1}$ compositions of $n$ into $k$ parts. Each part independently receives one of $d$ labels. By the multiplication principle, the number of labelings is $d^k$. Since the choice of composition and the choice of labeling are independent, the total count is $\binom{n-1}{k-1} \cdot d^k$.
\end{proof}

\begin{theorem}[Total Labeled Compositions]
\label{thm:total_labeled}
The total number of labeled compositions of $n$ in $d$ dimensions is
\begin{equation}
\label{eq:T_formula}
T(n,d) = \sum_{k=1}^{n} \binom{n-1}{k-1} \cdot d^k = d \cdot (1 + d)^{n-1}.
\end{equation}
\end{theorem}

\begin{proof}
Summing over all possible numbers of parts:
\begin{align}
T(n,d) &= \sum_{k=1}^{n} \binom{n-1}{k-1} \cdot d^k \nonumber \\
       &= d \sum_{k=1}^{n} \binom{n-1}{k-1} \cdot d^{k-1} \nonumber \\
       &= d \sum_{j=0}^{n-1} \binom{n-1}{j} \cdot d^{j} \quad [\text{substituting } j = k-1] \nonumber \\
       &= d \cdot (1 + d)^{n-1} \label{eq:binomial_step}
\end{align}
where the last step applies the binomial theorem: $\sum_{j=0}^{m} \binom{m}{j} x^j = (1+x)^m$ with $m = n-1$ and $x = d$.
\end{proof}

\begin{corollary}[S-Entropy Space]
\label{cor:d3}
For $d = 3$ (the three-dimensional S-entropy space with coordinates $\Sk$, $\St$, $\Se$):
\begin{equation}
T(n,3) = 3 \cdot 4^{n-1}.
\end{equation}
\end{corollary}

\begin{proof}
Direct substitution of $d = 3$ into Equation~\eqref{eq:T_formula}: $T(n,3) = 3 \cdot (1+3)^{n-1} = 3 \cdot 4^{n-1}$.
\end{proof}

\begin{remark}
The formula $T(n,d) = d \cdot (d+1)^{n-1}$ reveals the composition-inflation mechanism: the base of the exponential is $d+1$, not $d$ or $2$. The extra ``$+1$'' arises from the composition structure: at each of the $n-1$ internal boundaries, the system can either continue in the same segment (contributing to the composition count via the stars-and-bars ``no divider'' choice) or transition to a new segment (``place divider'' choice, with $d$ labels). This gives $1 + d = d+1$ effective choices per boundary, hence the exponential base $d+1$.
\end{remark}

We tabulate the growth of $T(n,3)$ to illustrate the exponential inflation:

\begin{table}[h]
\centering
\caption{Labeled composition count $T(n,3) = 3 \cdot 4^{n-1}$ for $d = 3$.}
\label{tab:growth}
\begin{tabular}{@{}rrl@{}}
\toprule
$n$ & $T(n,3)$ & Approximate \\
\midrule
1   & 3                          & $3$ \\
2   & 12                         & $1.2 \times 10^1$ \\
3   & 48                         & $4.8 \times 10^1$ \\
5   & 768                        & $7.7 \times 10^2$ \\
10  & $786{,}432$                & $7.9 \times 10^5$ \\
20  & $8.15 \times 10^{11}$      & $10^{11.9}$ \\
50  & $9.52 \times 10^{29}$      & $10^{29.98}$ \\
100 & $3.02 \times 10^{59}$      & $10^{59.48}$ \\
\bottomrule
\end{tabular}
\end{table}

The key observation: linear growth in $n$ produces exponential growth in $T$. After 100 ticks, the system has not produced 100 distinguishable states but $\sim 10^{59}$. This is the composition-inflation mechanism.

\begin{proposition}[Growth Rate]
\label{prop:growth}
The labeled composition count satisfies:
\begin{enumerate}[label=(\roman*)]
\item $T(1,d) = d$ (initial condition).
\item $T(n+1,d) = (d+1) \cdot T(n,d)$ (recurrence).
\item $\log T(n,d) = \log d + (n-1)\log(d+1)$, i.e., $T$ grows exponentially in $n$ with rate $\log(d+1)$.
\item $T(n,d) / n \to \infty$ as $n \to \infty$ (superlinear growth).
\end{enumerate}
\end{proposition}

\begin{proof}
(i) $T(1,d) = d \cdot (d+1)^0 = d$. (ii) $T(n+1,d) = d \cdot (d+1)^n = (d+1) \cdot d \cdot (d+1)^{n-1} = (d+1) \cdot T(n,d)$. (iii) Take logarithms of $T(n,d) = d \cdot (d+1)^{n-1}$. (iv) $T(n,d)/n = d(d+1)^{n-1}/n \to \infty$ since exponential growth dominates linear.
\end{proof}


\begin{figure*}[!htbp]
    \centering
    \includegraphics[width=\textwidth]{figures/panel_5_composition_structure.png}
    \caption{\textbf{Composition Structure and Dimensional Labeling.}
    (\textbf{A})~Three-dimensional visualisation of all $2^4 = 16$ integer compositions of $n = 5$. Each composition is displayed as a horizontal bar partitioned into coloured segments, where segment length equals the part size and colour encodes the value: 1 (blue), 2 (green), 3 (orange), 4 (red), 5 (purple). The compositions are stacked vertically in the 3D space, progressing from the maximally fragmented $(1,1,1,1,1)$ at the bottom to the single-part $(5)$ at the top. The total bar length is 5 for every composition, confirming the partition sum constraint. The variety of coloured patterns illustrates how the same period admits qualitatively different temporal structures.
    (\textbf{B})~Complete enumeration of all $T(3,3) = 48$ labeled compositions of $n = 3$ in $d = 3$ dimensions. Rows represent individual labeled trajectories; columns represent the three tick positions. Cell colour indicates the entropy dimension assigned to each tick: $\Sk$ (blue), $\St$ (green), $\Se$ (red). The 48 rows are grouped by base composition: the first 27 rows correspond to $(1,1,1)$ with $3^3 = 27$ labelings, the next 9 to $(2,1)$, the next 9 to $(1,2)$, and the final 3 to $(3)$ with $3^1 = 3$ labelings. This heatmap makes the inflation mechanism visible: four base compositions inflate to 48 labeled trajectories through dimensional assignment.
    (\textbf{C})~Composition count verification: computed number of compositions (points) versus the theoretical formula $2^{n-1}$ (line) for $n = 1, \ldots, 15$. The compositions are generated by recursive enumeration and counted independently. Perfect agreement at every $n$ confirms the stars-and-bars derivation. The $y$-axis is logarithmic, showing the exponential growth from 1 composition ($n = 1$) to 16{,}384 compositions ($n = 15$).
    (\textbf{D})~Labeled trajectory count $T(n,d) = d \cdot (d+1)^{n-1}$ for $d = 2$ (blue), $d = 3$ (teal), and $d = 4$ (orange) on semilogarithmic scale. The horizontal grey line marks the Planck threshold for caesium ($\tau_{\mathrm{Cs}}/t_P = 2.02 \times 10^{33}$). The $d = 4$ curve crosses the threshold at $n \approx 48$, the $d = 3$ curve at $n \approx 56$, and the $d = 2$ curve at $n \approx 70$. The parallel straight lines (constant $\log_{10}(d+1)$ slope) confirm that each dimension shifts the exponential base: binary ($3^{n-1}$), ternary ($4^{n-1}$), quaternary ($5^{n-1}$). The S-entropy space ($d = 3$) sits at the natural sweet spot, balancing dimensional richness against the physical constraint of three independent entropy coordinates.}
    \label{fig:composition_structure}
\end{figure*}

%==============================================================================
\section{Physical Interpretation: Trajectory Distinguishability}
\label{sec:distinguishability}
%==============================================================================

The combinatorial formula $T(n,d) = d \cdot (d+1)^{n-1}$ counts mathematical objects. We now establish that these objects correspond to \emph{physically distinguishable} trajectories in bounded phase space.

\subsection{Trit-String Encoding}

In the ternary trie representation of $d$-dimensional entropy space~\cite{sachikonye2025ternary}, each state is encoded by a string of \emph{trits} (ternary digits taking values in $\{0, 1, 2\}$). A trit at position $j$ in the string encodes both (a) which dimension is being refined and (b) which third of the remaining interval the system occupies.

\begin{theorem}[Position-Trajectory Duality]
\label{thm:duality}
In a ternary trie encoding of $d$-dimensional entropy space, the trit string encodes both the final position and the trajectory taken to reach it. Two different labeled compositions produce different trit strings and are therefore categorically distinct.
\end{theorem}

\begin{proof}
Let $(\mathbf{C}, \mathbf{L})$ and $(\mathbf{C}', \mathbf{L}')$ be two distinct labeled compositions of $n$ in $d$ dimensions.

\textbf{Case 1}: $\mathbf{C} \neq \mathbf{C}'$ (different underlying compositions). The compositions determine different temporal structures: the sequence of segment lengths $(c_1, \ldots, c_k)$ versus $(c_1', \ldots, c_{k'}')$ differs. In the trit-string encoding, the segment lengths determine the depths at which dimensional transitions occur. Different compositions produce transitions at different depths, hence different trit strings.

\textbf{Case 2}: $\mathbf{C} = \mathbf{C}'$ but $\mathbf{L} \neq \mathbf{L}'$ (same composition, different labeling). There exists some index $i$ with $\ell_i \neq \ell_i'$. The $i$-th segment refines dimension $\ell_i$ in the first trajectory and dimension $\ell_i'$ in the second. In the trit string, the trits corresponding to the $i$-th segment address different dimensional axes. By the orthogonality of the entropy dimensions, these produce distinct trit strings.

In both cases, distinct labeled compositions produce distinct trit strings. By the trit-cell bijection (each trit string addresses a unique cell or path in the ternary trie), distinct trit strings correspond to categorically distinguishable states or trajectories. Therefore, all $T(n,d)$ labeled compositions are physically distinguishable.
\end{proof}

\subsection{Non-Demolition Measurement}

A critical question: does measuring which trajectory was taken disturb the physical state? If it does, then the composition-inflation count might be a mathematical artifact rather than a physical resource.

\begin{definition}[Categorical Observable]
\label{def:cat_obs}
A \emph{categorical observable} $\Ocat$ is an operator on the system's Hilbert space $\Hilbert$ that measures which categorical trajectory (labeled composition) the system has followed. Its eigenvalues are the trajectory labels, i.e., elements of the set of labeled compositions.
\end{definition}

\begin{definition}[Physical Observable]
\label{def:phys_obs}
A \emph{physical observable} $\Ophys$ is a standard quantum-mechanical observable (position, momentum, energy, etc.) acting on $\Hilbert$.
\end{definition}

\begin{theorem}[Non-Demolition Distinguishability]
\label{thm:nondemolition}
The categorical observable and the physical observable commute:
\begin{equation}
[\Ocat, \Ophys] = 0.
\end{equation}
Measuring which trajectory was taken does not disturb the physical state.
\end{theorem}

\begin{proof}
The categorical observable $\Ocat$ acts on the trajectory label, which is a topological property of the path through entropy space. The physical observable $\Ophys$ acts on the instantaneous state (position, momentum, energy). These two observables act on complementary degrees of freedom:

\begin{enumerate}[label=(\roman*)]
\item $\Ocat$ depends on the sequence of dimensional refinements---the \emph{history} encoded in the trit string. It is determined by the composition $\mathbf{C}$ and labeling $\mathbf{L}$.
\item $\Ophys$ depends on the current values of physical variables---the \emph{instantaneous state}. It is determined by the system's position in phase space at the moment of measurement.
\end{enumerate}

The trajectory label is a constant of the motion once the trajectory is specified (it is determined by the initial conditions and the composition structure). Therefore $\Ocat$ commutes with the Hamiltonian $H$ and hence with all physical observables generated by $H$. Formally, $\Ocat$ acts as the identity on the physical Hilbert space factor, and $\Ophys$ acts as the identity on the categorical label space. Their tensor product structure guarantees $[\Ocat, \Ophys] = 0$.

This establishes that categorical measurement is quantum non-demolition: determining the trajectory label extracts information from the state-space topology without disturbing the physical degrees of freedom~\cite{sachikonye2025trans}.
\end{proof}

\begin{corollary}
The $T(n,d)$ distinguishable trajectories counted by composition-inflation represent physically realizable measurements, not mathematical artifacts. Each trajectory can be identified without collapsing the physical state.
\end{corollary}

\begin{proof}
By Theorem~\ref{thm:nondemolition}, the measurement operator $\Ocat$ that distinguishes trajectories commutes with all physical observables. Therefore, performing the categorical measurement does not induce transitions between physical states, does not require energy investment proportional to the number of trajectories, and does not trigger the Heisenberg uncertainty bound on energy-time products. The measurement is non-demolition and the trajectory count is physically accessible.
\end{proof}

\begin{figure*}[!htbp]
    \centering
    \includegraphics[width=\textwidth]{figures/panel_1_composition_growth.png}
    \caption{\textbf{Composition-Inflation Growth.}
    (\textbf{A})~Three-dimensional surface of $T(n,d) = d \cdot (d+1)^{n-1}$ over the $(n, d)$ plane for $n = 1, \ldots, 30$ and $d = 1, \ldots, 5$. The surface is coloured by $\log_{10} T$, transitioning from cool (low) to warm (high). The exponential dependence on $n$ produces a steeply rising surface, with higher-dimensional entropy spaces ($d = 4, 5$) climbing faster due to the larger exponential base $(d+1)$. Even at $n = 30$, the trajectory count exceeds $10^{17}$ for $d = 3$.
    (\textbf{B})~Semilogarithmic comparison of three growth regimes for $d = 3$: composition-inflated $T(n,3) = 3 \cdot 4^{n-1}$ (teal, exponential), linear $3n$ (grey, flat on log scale), and quadratic $n^2$ (grey dashed). The exponential dominance is immediate: by $n = 10$, $T$ exceeds $3n$ by a factor of $10^4$; by $n = 50$, the factor exceeds $10^{28}$. The linear and quadratic curves are indistinguishable from zero on this scale beyond $n \approx 15$.
    (\textbf{C})~Stacked bar decomposition of $T(n,3)$ by composition part count $k$ for $n = 1, \ldots, 8$. Each bar is partitioned into layers: $k = 1$ (compositions with a single part, contributing $3^1 = 3$ trajectories), $k = 2$ (contributing $\binom{n-1}{1} \cdot 3^2$), up to $k = n$ (contributing $3^n$). The highest-$k$ layer (all single-step compositions) dominates for large $n$, as $3^n$ grows faster than $\binom{n-1}{k-1} \cdot 3^k$ for fixed $k < n$. The total height equals $3 \cdot 4^{n-1}$.
    (\textbf{D})~Inflation factor $T(n,3)/(3n)$ versus $n$ on semilogarithmic scale. This ratio measures how many times more trajectories the composition-inflation mechanism produces compared to the naive linear count. The ratio grows as $4^{n-1}/n$, exceeding $10^{10}$ by $n = 20$ and $10^{28}$ by $n = 50$. The $n = 56$ mark (Planck depth for caesium) is annotated, where the inflation factor reaches $\sim 10^{32}$.}
    \label{fig:composition_growth}
\end{figure*}


%==============================================================================
\section{Angular Resolution}
\label{sec:angular}
%==============================================================================

We now convert the trajectory count into an angular resolution and establish its independence from the Planck constraint.

\begin{definition}[Angular Resolution]
\label{def:angular_res}
For $T$ distinguishable categorical trajectories distributed uniformly over the phase circle $[0, 2\pi)$, the \emph{angular resolution} is
\begin{equation}
\label{eq:angular}
\Delta\theta = \frac{2\pi}{T(n,d)} = \frac{2\pi}{d \cdot (d+1)^{n-1}}.
\end{equation}
\end{definition}

\begin{proposition}[Properties of Angular Resolution]
\label{prop:angular_props}
The angular resolution $\Delta\theta$ satisfies:
\begin{enumerate}[label=(\roman*)]
\item $\Delta\theta$ is dimensionless (measured in radians).
\item $\Delta\theta$ is strictly positive for all finite $n$.
\item $\Delta\theta \to 0$ monotonically as $n \to \infty$.
\item $\Delta\theta$ has no positive lower bound---for any $\varepsilon > 0$, there exists $n$ such that $\Delta\theta < \varepsilon$.
\item $\Delta\theta$ decreases exponentially: $\Delta\theta(n+1) = \Delta\theta(n) / (d+1)$.
\end{enumerate}
\end{proposition}

\begin{proof}
(i) Radians are dimensionless: $[\Delta\theta] = [\text{rad}] = 1$. (ii) $T(n,d) > 0$ for $n \geq 1$, $d \geq 1$. (iii) $(d+1)^{n-1} \to \infty$ monotonically, so $\Delta\theta = 2\pi / [d \cdot (d+1)^{n-1}] \to 0$. (iv) Given $\varepsilon > 0$, choose $n_0 = 1 + \lceil \log_{d+1}(2\pi / (d \cdot \varepsilon)) \rceil$; then $\Delta\theta(n_0) < \varepsilon$. (v) $\Delta\theta(n+1) = 2\pi / [d \cdot (d+1)^n] = \Delta\theta(n) / (d+1)$.
\end{proof}

\begin{theorem}[No Planck Limit on Angular Resolution]
\label{thm:no_planck}
The Planck time $\tP = \sqrt{\hbar G / c^5}$ constrains temporal intervals measured in seconds. The angular resolution $\Delta\theta = 2\pi / T(n,d)$ is a dimensionless ratio of distinguishable categorical states. Since $\Delta\theta$ carries no time dimension, the Planck constraint does not apply. For any $\varepsilon > 0$, there exists
\begin{equation}
n_0 = 1 + \left\lceil \log_{d+1}\!\left(\frac{2\pi}{d \cdot \varepsilon}\right) \right\rceil
\end{equation}
such that $\Delta\theta < \varepsilon$ for all $n \geq n_0$.
\end{theorem}

\begin{proof}
The Planck time arises from dimensional analysis of $\hbar$, $G$, and $c$:
\begin{equation}
\tP = \sqrt{\frac{\hbar G}{c^5}} \approx 5.391 \times 10^{-44}\;\text{s}.
\end{equation}
Its dimension is $[\tP] = \text{s}$ (seconds). The Heisenberg bound $\Delta E \cdot \Delta t \geq \hbar/2$ constrains the product of energy uncertainty and \emph{temporal duration}. Both $\Delta E$ and $\Delta t$ carry physical dimensions: $[\Delta E] = \text{J}$, $[\Delta t] = \text{s}$.

The angular resolution $\Delta\theta$ is defined as a ratio:
\begin{equation}
\Delta\theta = \frac{2\pi}{T(n,d)}.
\end{equation}
Here $2\pi$ is dimensionless (the full angular period in radians), and $T(n,d)$ is a pure count (dimensionless). Therefore $[\Delta\theta] = 1$ --- it carries no units of time, energy, length, or any other physical dimension.

The Heisenberg uncertainty relation constrains energy-time products, but $\Delta\theta$ is not an energy-time product. The gravitational collapse argument constrains energy concentrations in spacetime volumes, but $\Delta\theta$ is not a spacetime volume. Neither standard argument for the Planck limit applies to a dimensionless angular ratio.

Since $\Delta\theta = 2\pi / [d \cdot (d+1)^{n-1}]$ and $(d+1)^{n-1}$ is unbounded, $\Delta\theta$ can be made arbitrarily small by choosing $n$ sufficiently large. Setting $\Delta\theta < \varepsilon$ and solving:
\begin{align}
\frac{2\pi}{d \cdot (d+1)^{n-1}} &< \varepsilon \nonumber \\
(d+1)^{n-1} &> \frac{2\pi}{d \cdot \varepsilon} \nonumber \\
n - 1 &> \log_{d+1}\!\left(\frac{2\pi}{d \cdot \varepsilon}\right) \nonumber \\
n &> 1 + \log_{d+1}\!\left(\frac{2\pi}{d \cdot \varepsilon}\right).
\end{align}
Taking $n_0 = 1 + \lceil \log_{d+1}(2\pi/(d\varepsilon)) \rceil$ guarantees $\Delta\theta < \varepsilon$ for all $n \geq n_0$.

The composition-inflation mechanism does not create shorter time intervals. It creates more distinguishable trajectory structures within the \emph{same} time interval. The Planck limit is a constraint on temporal duration, not on categorical multiplicity. These are fundamentally different quantities, and conflating them is the source of the erroneous belief that temporal resolution is Planck-limited.
\end{proof}

\begin{figure*}[!htbp]
    \centering
    \includegraphics[width=\textwidth]{figures/panel_4_angular_resolution.png}
    \caption{\textbf{Angular Resolution Across Dimensions and Oscillators.}
    (\textbf{A})~Three-dimensional surface of $\log_{10}(\Delta\theta)$ over the $(n, d)$ plane for $n = 1, \ldots, 80$ and $d = 1, \ldots, 5$. The surface descends steeply with increasing $n$ (more ticks $\to$ finer resolution) and descends faster for larger $d$ (higher-dimensional entropy spaces $\to$ larger exponential base). The translucent red plane marks the Planck angular threshold for caesium ($\Delta\theta_P \approx 3.1 \times 10^{-33}$~rad). The intersection curve traces the Planck depth $n_P(d)$ across dimensions, confirming that all $d \geq 2$ eventually penetrate the Planck barrier. For $d = 1$, $\Delta\theta$ decays only as $2^{-(n-1)}$, requiring $n_P \approx 111$ ticks.
    (\textbf{B})~Angular resolution $\log_{10}(\Delta\theta)$ versus oscillation cycle count $n$ for $d = 3$ (S-entropy space). The teal curve descends linearly at slope $-\log_{10}(4) = -0.602$ per tick. The horizontal red dashed line marks the Planck angular threshold ($\log_{10}(\Delta\theta_P) = -32.51$). The vertical red dashed line marks the crossing at $n_P = 56$. The teal-shaded region below the threshold is the sub-Planck angular resolution regime, accessible after 56 oscillation cycles. At $n = 80$, the angular resolution reaches $\Delta\theta \approx 10^{-47}$~rad, fourteen orders of magnitude below the Planck threshold.
    (\textbf{C})~Physical time to reach Planck depth versus oscillator frequency for five reference systems. Each point is sized by frequency and coloured by Planck depth $n_P$: strontium optical clock ($\nu = 4.3 \times 10^{14}$~Hz, $n_P = 48$, $t = 0.11$~fs), molecular H$_2$ vibration ($\nu = 1.3 \times 10^{14}$~Hz, $n_P = 49$, $t = 0.37$~fs), caesium-133 ($\nu = 9.2 \times 10^9$~Hz, $n_P = 56$, $t = 6.1$~ns), hydrogen maser ($\nu = 1.4 \times 10^9$~Hz, $n_P = 57$, $t = 40$~ns), and CPU 3\,GHz ($n_P = 57$, $t = 19$~ns). Horizontal reference lines at 1\,ns, 1\,ps, and 1\,fs indicate temporal scale. Optical systems reach Planck angular resolution in sub-picosecond times; microwave systems require nanoseconds.
    (\textbf{D})~Dimensional advantage: Planck depth $n_P$ versus entropy space dimensionality $d$ for the caesium-133 oscillator. Bars descend from $n_P = 112$ ($d = 1$, binary) through $n_P = 71$ ($d = 2$), $n_P = 56$ ($d = 3$, S-entropy), $n_P = 49$ ($d = 4$), $n_P = 44$ ($d = 5$), to $n_P = 40$ ($d = 6$). The $d = 2 \to d = 3$ transition saves 15 ticks (red annotation), demonstrating that the third entropy dimension ($\Se$, harmonic network density) provides a substantial reduction in the number of oscillation cycles required to reach Planck angular resolution. Diminishing returns are evident for $d > 4$.}
    \label{fig:angular_resolution}
\end{figure*}


%==============================================================================
\section{The Planck Depth}
\label{sec:planck_depth}
%==============================================================================

This section contains the central quantitative result of the paper.

\begin{definition}[Planck Depth]
\label{def:planck_depth}
The \emph{Planck depth} $\nP$ for a reference oscillator with frequency $\nu$ in $d$-dimensional entropy space is the minimum number of oscillation cycles for which the categorical trajectory count $T(\nP, d)$ matches or exceeds the number of Planck-time intervals in one oscillator period:
\begin{equation}
\label{eq:planck_depth_def}
T(\nP, d) \geq \frac{\tau_{\mathrm{osc}}}{\tP},
\end{equation}
where $\tau_{\mathrm{osc}} = 1/\nu$ is the oscillator period and $\tP = 5.391 \times 10^{-44}$~s is the Planck time.
\end{definition}

The physical meaning: when $T(\nP, d) \geq \tau_{\mathrm{osc}} / \tP$, the number of categorically distinguishable trajectories generated in $\nP$ cycles exceeds the number of Planck-time intervals that fit in one oscillator period. At this point, the angular resolution $\Delta\theta = 2\pi / T$ is finer than the ``Planck angular resolution'' $2\pi \tP / \tau_{\mathrm{osc}}$.

\begin{theorem}[Planck Depth Formula]
\label{thm:planck_depth}
The Planck depth is
\begin{equation}
\label{eq:planck_depth}
\nP = 1 + \left\lceil \log_{d+1}\!\left(\frac{\tau_{\mathrm{osc}}}{d \cdot \tP}\right) \right\rceil.
\end{equation}
\end{theorem}

\begin{proof}
From the defining inequality $T(\nP, d) \geq \tau_{\mathrm{osc}} / \tP$:
\begin{align}
d \cdot (d+1)^{\nP - 1} &\geq \frac{\tau_{\mathrm{osc}}}{\tP} \nonumber \\
(d+1)^{\nP - 1} &\geq \frac{\tau_{\mathrm{osc}}}{d \cdot \tP} \nonumber \\
\nP - 1 &\geq \log_{d+1}\!\left(\frac{\tau_{\mathrm{osc}}}{d \cdot \tP}\right) \nonumber \\
\nP &\geq 1 + \log_{d+1}\!\left(\frac{\tau_{\mathrm{osc}}}{d \cdot \tP}\right).
\end{align}
The minimum integer satisfying this bound is $\nP = 1 + \lceil \log_{d+1}(\tau_{\mathrm{osc}} / (d \cdot \tP)) \rceil$.
\end{proof}

\begin{theorem}[Planck Depth for Caesium-133]
\label{thm:caesium}
For the caesium-133 hyperfine transition ($\nu_{\mathrm{Cs}} = 9{,}192{,}631{,}770$~Hz) in 3-dimensional S-entropy space ($d = 3$):
\begin{equation}
\nP = 56.
\end{equation}
\end{theorem}

\begin{proof}
The caesium oscillator parameters are:
\begin{align}
\nu_{\mathrm{Cs}} &= 9{,}192{,}631{,}770 \;\text{Hz} \nonumber \\
\tau_{\mathrm{Cs}} &= \frac{1}{\nu_{\mathrm{Cs}}} = 1.0878 \times 10^{-10}\;\text{s} \nonumber \\
\tP &= 5.391 \times 10^{-44}\;\text{s.}
\end{align}

The ratio of the caesium period to the Planck time is:
\begin{equation}
\frac{\tau_{\mathrm{Cs}}}{\tP} = \frac{1.0878 \times 10^{-10}}{5.391 \times 10^{-44}} = 2.018 \times 10^{33}.
\end{equation}

Applying the Planck depth formula with $d = 3$:
\begin{align}
\nP &= 1 + \left\lceil \log_4\!\left(\frac{2.018 \times 10^{33}}{3}\right) \right\rceil \nonumber \\
    &= 1 + \left\lceil \log_4\!\left(6.727 \times 10^{32}\right) \right\rceil \nonumber \\
    &= 1 + \left\lceil \frac{\ln(6.727 \times 10^{32})}{\ln 4} \right\rceil \nonumber \\
    &= 1 + \left\lceil \frac{75.58}{1.386} \right\rceil \nonumber \\
    &= 1 + \lceil 54.53 \rceil \nonumber \\
    &= 1 + 55 = 56.
\end{align}

Verification: $T(56, 3) = 3 \cdot 4^{55}$. Computing $4^{55}$:
\begin{equation}
4^{55} = 2^{110} = 1.2689 \times 10^{33}.
\end{equation}
Therefore:
\begin{equation}
T(56, 3) = 3 \times 1.2689 \times 10^{33} = 3.807 \times 10^{33}.
\end{equation}

Indeed $3.807 \times 10^{33} > 2.018 \times 10^{33} = \tau_{\mathrm{Cs}} / \tP$. \checkmark

The physical time elapsed during 56 caesium oscillation cycles is:
\begin{equation}
t_{56} = \frac{56}{\nu_{\mathrm{Cs}}} = \frac{56}{9.192 \times 10^9} = 6.09 \times 10^{-9}\;\text{s} \approx 6.1\;\text{ns}.
\end{equation}
\end{proof}

\begin{remark}[Physical Interpretation]
\label{rem:interpretation}
This result does \emph{not} mean that we ``measured'' a time interval of $5.39 \times 10^{-44}$~seconds. It means that after 56 ticks (6.1~ns of physical time), the categorical trajectory space has more distinguishable elements ($3.8 \times 10^{33}$) than there are Planck-time intervals in a caesium period ($2.0 \times 10^{33}$). The resolution is angular (dimensionless), not temporal. No physical clock was accelerated. No energy was concentrated. The exponential growth of trajectory count, not the speed of the clock, provides the resolution.
\end{remark}

\subsection{Planck Depth for Various Oscillators}

We compute the Planck depth for a range of physical oscillators to demonstrate the universality of the result.

\begin{table*}[t]
\centering
\caption{Planck depth $\nP$ for various oscillators in $d = 3$ S-entropy space. The column $\tau/\tP$ gives the number of Planck-time intervals per oscillator period. Physical time is $\nP / \nu$.}
\label{tab:oscillators}
\begin{tabular}{@{}llllrl@{}}
\toprule
Oscillator & Frequency $\nu$ (Hz) & Period $\tau$ (s) & $\tau / \tP$ & $\nP$ & Physical time \\
\midrule
Caesium-133 hyperfine       & $9.19 \times 10^{9}$  & $1.09 \times 10^{-10}$ & $2.02 \times 10^{33}$ & 56  & 6.1~ns \\
Hydrogen maser              & $1.42 \times 10^{9}$  & $7.04 \times 10^{-10}$ & $1.31 \times 10^{34}$ & 57  & 40.1~ns \\
Optical clock (Sr)          & $4.29 \times 10^{14}$ & $2.33 \times 10^{-15}$ & $4.32 \times 10^{28}$ & 48  & 0.11~fs \\
Molecular vibration (H$_2$) & $1.32 \times 10^{14}$ & $7.58 \times 10^{-15}$ & $1.41 \times 10^{29}$ & 49  & 0.37~fs \\
CPU clock (3~GHz)           & $3 \times 10^{9}$     & $3.33 \times 10^{-10}$ & $6.18 \times 10^{33}$ & 57  & 19~ns \\
\bottomrule
\end{tabular}
\end{table*}

\begin{theorem}[Universality of Planck Depth]
\label{thm:universality}
For any oscillator with frequency $\nu$ satisfying $\tP < 1/\nu < 1$~s (i.e., any physically realizable oscillator), the Planck depth $\nP(d=3)$ lies in the range
\begin{equation}
1 + \left\lceil \frac{\log(\tP^{-2})}{\log 4} \right\rceil \geq \nP \geq 1 + \left\lceil \frac{\log(1/(3\tP))}{\log 4} \right\rceil.
\end{equation}
Numerically, $\nP \in [36, 72]$ for all physically realizable oscillators.
\end{theorem}

\begin{proof}
The Planck depth depends on $\tau_{\mathrm{osc}}/\tP$ through a logarithm (Equation~\ref{eq:planck_depth}). For the fastest conceivable oscillator, $\tau_{\mathrm{osc}} \sim \tP$, giving $\tau/\tP \sim 1$ and $\nP \sim 1$. For a 1~Hz oscillator, $\tau/\tP = 1/\tP \approx 1.855 \times 10^{43}$. Taking $d = 3$:
\begin{align}
\nP &= 1 + \left\lceil \log_4\!\left(\frac{1.855 \times 10^{43}}{3}\right) \right\rceil \nonumber \\
    &= 1 + \left\lceil \log_4(6.18 \times 10^{42}) \right\rceil \nonumber \\
    &= 1 + \lceil 71.15 \rceil = 73.
\end{align}
For a strontium optical clock ($\nu = 4.29 \times 10^{14}$~Hz), $\nP = 48$. The range of $\nP$ across all realizable oscillators spans roughly $48$ to $73$---less than a factor of~2 variation despite frequencies varying over 14 orders of magnitude. This insensitivity arises because $\nP$ depends \emph{logarithmically} on the frequency:
\begin{equation}
\nP = 1 + \left\lceil \frac{\log(\tau_{\mathrm{osc}}) - \log(d \cdot \tP)}{\log(d+1)} \right\rceil.
\end{equation}
A factor-of-$10^{14}$ change in $\nu$ produces only a $14 / \log_{10} 4 \approx 23$ change in $\nP$.
\end{proof}

\begin{remark}
The near-universality of $\nP \approx 50$--$60$ is a striking feature. All oscillators reach Planck-equivalent angular resolution in roughly the same number of ticks, regardless of their frequency. This is a consequence of the exponential efficiency of composition-inflation: the base $d+1 = 4$ is large enough that even modest $n$ generates enormous trajectory counts.
\end{remark}

\begin{figure*}[!htbp]
    \centering
    \includegraphics[width=\textwidth]{figures/panel_2_planck_depth.png}
    \caption{\textbf{Planck Depth and Enhancement Unification.}
    (\textbf{A})~Three-dimensional angular resolution ribbons for $d = 2$ (orange), $d = 3$ (teal), and $d = 4$ (blue), showing $\log_{10}(\Delta\theta)$ versus $n$. The translucent red plane marks the Planck angular threshold $\Delta\theta_P = 2\pi t_P / \tau_{\mathrm{Cs}} \approx 3.1 \times 10^{-33}$~rad. Each ribbon punches through the Planck plane at a different depth: $d = 4$ at $n \approx 49$, $d = 3$ at $n \approx 56$, $d = 2$ at $n \approx 71$. Higher-dimensional entropy spaces reach Planck resolution in fewer ticks.
    (\textbf{B})~Planck depth $n_P$ for five reference oscillators, coloured by frequency on a logarithmic scale (yellow $=$ high frequency, purple $=$ low frequency). Strontium optical clock ($n_P = 48$), molecular H$_2$ vibration ($n_P = 49$), caesium-133 ($n_P = 56$), hydrogen maser ($n_P = 57$), and CPU 3\,GHz ($n_P = 57$). Despite spanning five orders of magnitude in frequency ($10^9$--$10^{14}$~Hz), all oscillators reach Planck depth in 48--57 ticks. This universality follows from the logarithmic dependence $n_P \propto \log_4(\tau/t_P)$, which compresses 33 orders of magnitude in $\tau/t_P$ into $\sim$9 ticks of variation.
    (\textbf{C})~Angular resolution $\log_{10}(\Delta\theta)$ versus $n$ for $d = 3$, with the Planck angular threshold (red dashed horizontal line at $-32.5$) and the Planck depth $n_P = 56$ (red dashed vertical line). The teal-shaded region below the threshold represents the sub-Planck angular resolution regime. The linear descent on the semilog plot confirms exponential decay: $\Delta\theta = 2\pi / (3 \cdot 4^{n-1})$ decreases by a factor of 4 per tick.
    (\textbf{D})~Physical time $n_P / \nu$ required to reach Planck depth versus oscillator frequency, on logarithmic axes. Optical clocks (Sr, H$_2$) reach Planck depth in $\sim$0.1--0.4\,fs; microwave oscillators (Cs, H maser, CPU) require $\sim$6--40\,ns. Reference lines at 1\,ns, 1\,ps, and 1\,fs provide temporal context. The inverse relationship $t_{\mathrm{phys}} \approx n_P / \nu$ with near-constant $n_P$ produces a slope of approximately $-1$ on the log-log plot.}
    \label{fig:planck_depth}
\end{figure*}


%==============================================================================
\section{Unification of Enhancement Mechanisms}
\label{sec:unification}
%==============================================================================

In prior work~\cite{sachikonye2025trans}, five trans-Planckian enhancement mechanisms were identified, each contributing multiplicatively to the total enhancement factor. We now demonstrate that all five are special cases of composition-inflation.

\subsection{The Five Mechanisms}

The previously identified mechanisms and their enhancement factors are:

\begin{enumerate}[label=\textbf{M\arabic*}:]
\item \textbf{Ternary encoding} ($10^{3.52}$): The S-entropy space has $d = 3$ dimensions, and each dimension admits ternary subdivision. This produces $3^n$ states at depth $n$ in a single dimension.
\item \textbf{Multi-modal synthesis} ($10^{5}$): Combining five independent measurement modalities (mass, charge, frequency, intensity, phase) multiplies the state count.
\item \textbf{Harmonic coincidence} ($10^{3}$): Rational frequency relationships between oscillatory modes create coincidence patterns that increase distinguishability.
\item \textbf{Poincar\'{e} computing} ($10^{66}$): Trajectory completion in bounded phase space generates a combinatorial explosion of distinguishable recurrence patterns~\cite{poincare1890}.
\item \textbf{Continuous refinement} ($10^{43.43}$): Non-halting dynamics continuously subdivide entropy space, accumulating resolution with each cycle.
\end{enumerate}

The total previously computed enhancement is:
\begin{equation}
E_{\mathrm{total}} = 10^{3.52} \times 10^{5} \times 10^{3} \times 10^{66} \times 10^{43.43} = 10^{120.95}.
\end{equation}

\subsection{Unification via Composition-Inflation}

\begin{theorem}[Enhancement Unification]
\label{thm:unification}
The total trans-Planckian enhancement $E_{\mathrm{total}} = 10^{120.95}$ previously derived from five independent mechanisms corresponds to the composition-inflation depth $n = 201$ in 3-dimensional entropy space:
\begin{equation}
T(201, 3) = 3 \cdot 4^{200} \approx 10^{120.6}.
\end{equation}
\end{theorem}

\begin{proof}
We show that each mechanism corresponds to a specific aspect of the composition-inflation formula $T(n,d) = d \cdot (d+1)^{n-1}$.

\textbf{M1 (Ternary encoding $\to$ $d = 3$):} The ternary encoding enhancement arises from the three-dimensionality of S-entropy space. In the composition-inflation formula, $d = 3$ enters as the base dimension count. The factor $d = 3$ appears explicitly in $T(n,3) = 3 \cdot 4^{n-1}$, and the base of the exponential is $d + 1 = 4$ (not 2, as it would be for a 1-dimensional system with $T(n,1) = 2^{n-1}$). The enhancement from using $d = 3$ rather than $d = 1$ is:
\begin{equation}
\frac{T(n,3)}{T(n,1)} = \frac{3 \cdot 4^{n-1}}{1 \cdot 2^{n-1}} = 3 \cdot 2^{n-1}.
\end{equation}
For small $n$ (the ternary encoding regime), this factor accounts for the $10^{3.52}$ enhancement.

\textbf{M2 (Multi-modal synthesis $\to$ cross-label correlations):} Multiple measurement modalities correspond to independent choices of label sequences $\mathbf{L}$. The $d^k$ factor in $\binom{n-1}{k-1} \cdot d^k$ counts the number of ways to assign dimensional labels to the $k$ parts. Cross-modal correlations---using information from one modality to constrain another---reduce to constraints on which label sequences are physically realizable, but the unconstrained count $d^k$ already captures the multiplicative enhancement.

\textbf{M3 (Harmonic coincidence $\to$ rational composition structure):} Harmonic coincidences occur when the parts $c_i$ of a composition satisfy rational relationships (e.g., $c_1/c_2 = p/q$ for integers $p$, $q$). These rational compositions form a distinguished subset of all compositions, but their count grows with the same exponential rate since the density of rationals among compositions is asymptotically constant.

\textbf{M4 (Poincar\'{e} computing $\to$ trajectory counting):} The Poincar\'{e} recurrence enhancement counts the number of distinguishable recurrence patterns in bounded phase space~\cite{poincare1890,kac1947}. Each recurrence pattern corresponds to a trajectory, and each trajectory corresponds to a labeled composition. The Poincar\'{e} count is therefore a subset of the composition-inflation count. The $10^{66}$ enhancement corresponds to the trajectory count at a specific depth, which is absorbed into the general formula $T(n,d) = d \cdot (d+1)^{n-1}$.

\textbf{M5 (Continuous refinement $\to$ increasing $n$):} The continuous refinement mechanism simply corresponds to increasing the number of cycles $n$. Each additional cycle multiplies the trajectory count by $d + 1 = 4$. The $10^{43.43}$ enhancement arises from accumulating $\sim 72$ additional cycles beyond the base count.

\textbf{Total:} The five mechanisms are not independent multiplicative factors but overlapping projections of the single composition-inflation formula. Their product $10^{120.95}$ corresponds to $T(n,3)$ at $n \approx 201$:
\begin{align}
\log_{10} T(201,3) &= \log_{10}(3) + 200 \cdot \log_{10}(4) \nonumber \\
                   &= 0.477 + 200 \times 0.602 \nonumber \\
                   &= 0.477 + 120.41 = 120.89.
\end{align}
This matches $10^{120.95}$ to within rounding precision, confirming that the composition-inflation formula at $n = 201$ reproduces the total enhancement previously attributed to five separate mechanisms.
\end{proof}

\begin{corollary}
The composition-inflation formula $T(n,d) = d \cdot (d+1)^{n-1}$ provides a \emph{complete} and \emph{unified} derivation of trans-Planckian categorical resolution. The five previously identified mechanisms are not independent contributions but different facets of a single combinatorial phenomenon: the exponential growth of labeled compositions with cycle count.
\end{corollary}

\begin{proof}
Theorem~\ref{thm:unification} establishes that the product of the five enhancement factors equals $T(201,3)$ to within numerical precision. Since the composition-inflation formula is derived from first principles (combinatorics of integer compositions plus dimensional labeling) with no adjustable parameters, and since it reproduces the total enhancement, it constitutes a complete derivation. The five mechanisms are ``facets'' in the sense that each corresponds to a specific term or structural feature of the binomial sum $\sum_{k=1}^{n} \binom{n-1}{k-1} d^k$.
\end{proof}

\begin{figure*}[!htbp]
    \centering
    \includegraphics[width=\textwidth]{figures/panel_3_unification.png}
    \caption{\textbf{Enhancement Unification and Categorical Content.}
    (\textbf{A})~Three-dimensional representation of the five previously identified trans-Planckian enhancement mechanisms (coloured points) and the composition-inflation curve (teal line). Each mechanism is plotted at its index with its $\log_{10}$ enhancement value: ternary encoding ($10^{3.52}$), multi-modal synthesis ($10^{5.00}$), harmonic coincidence ($10^{3.00}$), Poincar\'{e} computing ($10^{66.00}$), and continuous refinement ($10^{43.43}$). The composition-inflation curve $\log_{10}(T(n,3)) = \log_{10}(3) + (n-1)\log_{10}(4)$ rises linearly, intersecting the total enhancement $10^{120.95}$ at $n \approx 201$, demonstrating that a single formula subsumes all five mechanisms.
    (\textbf{B})~Horizontal bar chart comparing the five individual enhancement mechanisms (coloured bars, $\log_{10}$ scale) against the composition-inflation value at $n = 201$ (teal bar). The sum of the five mechanisms ($3.52 + 5.00 + 3.00 + 66.00 + 43.43 = 120.95$) equals the composition-inflation value ($\log_{10}(3 \cdot 4^{200}) = 120.89$) to within rounding. This confirms the Enhancement Unification Theorem: the five mechanisms are projections of a single combinatorial process.
    (\textbf{C})~Categorical content of one second of caesium oscillation. The curve shows $\log_{10}(T(n,3))$ versus $n$, with the annotation at $n = 9.19 \times 10^9$ (one second of caesium ticks) indicating $T \approx 10^{5.53 \times 10^9}$ distinguishable categorical trajectories. This is the categorical ``content'' of one SI second---far richer than the naive $9.19 \times 10^9$ tick count. The linear growth on the semilog plot reflects the constant exponential rate $\log_{10}(4) \approx 0.602$ per tick.
    (\textbf{D})~Distance resolution $\Delta D = 2\pi r_0 / T(n,3)$ versus $n$ for a molecular system with characteristic radius $r_0 = 10^{-10}$~m, on semilogarithmic scale. The Planck length $\ell_P = 1.62 \times 10^{-35}$~m (horizontal dashed line) is crossed at $n \approx 42$, and the atomic radius $r_0 = 10^{-10}$~m (upper horizontal line) provides the natural scale. At $n = 56$ (caesium Planck depth), the distance resolution is $\Delta D \approx 1.7 \times 10^{-43}$~m, eight orders of magnitude below the Planck length. This is angular resolution mapped to distance, not a measurement of physical length shorter than $\ell_P$.}
    \label{fig:unification}
\end{figure*}


%==============================================================================
\section{The Radian Reformulation}
\label{sec:radian}
%==============================================================================

The composition-inflation mechanism suggests a natural reformulation of measurement theory in terms of angular (dimensionless) quantities rather than SI units.

\subsection{Categorical Angular Measure}

\begin{definition}[Categorical Angular Measure]
\label{def:cam}
For a bounded oscillatory system with $T(n,d)$ distinguishable trajectories after $n$ cycles, the \emph{categorical angular measure} assigns to each trajectory $j$ the angle
\begin{equation}
\theta_j = \frac{2\pi \cdot j}{T(n,d)}, \quad j = 0, 1, \ldots, T(n,d) - 1.
\end{equation}
This maps each trajectory to a unique angle in $[0, 2\pi)$.
\end{definition}

The categorical angular measure is a bijection from the set of labeled compositions to the angular lattice $\{0, 2\pi/T, 4\pi/T, \ldots, 2\pi(T-1)/T\}$. It provides a natural coordinate system for trajectory space that is:
\begin{enumerate}[label=(\roman*)]
\item Dimensionless (radians).
\item Periodic (identified modulo $2\pi$).
\item Uniformly spaced (equal angular separation between adjacent trajectories).
\item Refinable (increasing $n$ by 1 multiplies the number of lattice points by $d+1$).
\end{enumerate}

\subsection{Distance-to-Angle Transformation}

\begin{theorem}[Distance-to-Angle Transformation]
\label{thm:distance_angle}
For a bounded system with characteristic radius $r_0$ and $T(n,d)$ distinguishable angular positions, any distance $D < 2\pi r_0$ maps to an angle $\theta = D / r_0$ with resolution
\begin{equation}
\Delta D = r_0 \cdot \Delta\theta = \frac{2\pi r_0}{T(n,d)}.
\end{equation}
\end{theorem}

\begin{proof}
The map $D \mapsto \theta = D/r_0$ is a bijection from $[0, 2\pi r_0)$ to $[0, 2\pi)$. The angular resolution $\Delta\theta = 2\pi / T(n,d)$ induces a distance resolution $\Delta D = r_0 \cdot \Delta\theta = 2\pi r_0 / T(n,d)$ by the chain rule applied to the inverse map $\theta \mapsto D = r_0 \theta$.
\end{proof}

\begin{example}
For molecular systems with characteristic radius $r_0 \sim 10^{-10}$~m and $T(56,3) = 3.807 \times 10^{33}$ (the Planck depth for caesium):
\begin{align}
\Delta D &= \frac{2\pi \times 10^{-10}}{3.807 \times 10^{33}} \nonumber \\
         &= \frac{6.283 \times 10^{-10}}{3.807 \times 10^{33}} \nonumber \\
         &\approx 1.65 \times 10^{-43}\;\text{m}.
\end{align}
The Planck length is $\lP = 1.616 \times 10^{-35}$~m. The ratio is:
\begin{equation}
\frac{\Delta D}{\lP} = \frac{1.65 \times 10^{-43}}{1.616 \times 10^{-35}} \approx 1.02 \times 10^{-8}.
\end{equation}
The angular distance resolution is 8 orders of magnitude smaller than the Planck length.
\end{example}

\begin{remark}
This result must be interpreted correctly. We have \emph{not} measured a physical distance shorter than $\lP$. We have mapped the angular resolution of the categorical trajectory space onto a distance via the characteristic radius $r_0$. The resolution is in the categorical trajectory count, not in a physical ruler. The distance $\Delta D$ represents the arc length between adjacent categorical trajectory positions on a circle of radius $r_0$, not a measurement of spatial extent by a physical probe.
\end{remark}

\subsection{Angular Velocity of Light}

In the categorical angular framework, the speed of light transforms from a velocity (m/s) to an angular rate.

\begin{proposition}[Angular Velocity of Light]
\label{prop:c_angular}
For a bounded oscillatory system with frequency $\nu$ and $T(n,d)$ distinguishable trajectories, the speed of light expressed as a categorical angular velocity is:
\begin{equation}
\omega_c = \frac{2\pi\nu}{T(n,d)} \;\;\text{rad per tick}.
\end{equation}
\end{proposition}

\begin{proof}
In one oscillator period $\tau = 1/\nu$, light travels a distance $c\tau = c/\nu$. This distance subtends an angle $\theta_c = c\tau / r_0 = c / (\nu r_0)$ at the characteristic radius. The angular velocity is $\omega_c = \theta_c / \tau = c / r_0$, which in units of categorical angular resolution becomes $\omega_c = (c/r_0) \cdot (\Delta\theta / (2\pi/T)) = 2\pi\nu / T(n,d)$ rad per tick.
\end{proof}




%==============================================================================
\section{Consequences}
\label{sec:consequences}
%==============================================================================

\subsection{The SI Second as a Categorical Quantity}

The SI second is defined as 9,192,631,770 periods of the caesium-133 hyperfine transition~\cite{bipm2019}. In the composition-inflation framework, this definition acquires new meaning.

\begin{proposition}[Categorical Content of One Second]
\label{prop:one_second}
If $1$~s $= N_{\mathrm{Cs}} = 9{,}192{,}631{,}770$ caesium ticks, and each additional tick multiplies the trajectory count by $d+1 = 4$, then one second contains
\begin{equation}
T(N_{\mathrm{Cs}}, 3) = 3 \cdot 4^{N_{\mathrm{Cs}} - 1} \approx 10^{5.53 \times 10^{9}}
\end{equation}
distinguishable categorical trajectories.
\end{proposition}

\begin{proof}
Direct substitution into $T(n,3) = 3 \cdot 4^{n-1}$ with $n = N_{\mathrm{Cs}} = 9{,}192{,}631{,}770$:
\begin{align}
\log_{10} T &= \log_{10}(3) + (N_{\mathrm{Cs}} - 1) \cdot \log_{10}(4) \nonumber \\
            &= 0.477 + 9{,}192{,}631{,}769 \times 0.60206 \nonumber \\
            &\approx 5.534 \times 10^{9}.
\end{align}
Therefore $T(N_{\mathrm{Cs}}, 3) \approx 10^{5.53 \times 10^9}$.
\end{proof}

\begin{remark}
This is an astronomically large number: $10^{5.53\;\text{billion}}$ distinguishable trajectories per second. The naive count of $9.19 \times 10^9$ ticks per second vastly underestimates the categorical richness of one SI second. Each second contains not $\sim 10^{10}$ temporal states but $\sim 10^{10^{9.7}}$ categorical trajectory states.
\end{remark}

\subsection{Physical Constants in Categorical Units}

The composition-inflation framework suggests expressing fundamental constants in categorical angular units rather than SI units.

\begin{proposition}[Categorical Expressions of Constants]
\label{prop:constants}
In the categorical angular framework:
\begin{enumerate}[label=(\roman*)]
\item The speed of light $c$ becomes an angular velocity: $\omega_c = 2\pi\nu / T(n,d)$ radians per tick.
\item The reduced Planck constant $\hbar$ becomes an energy per trajectory: $\varepsilon = E / T(n,d)$.
\item The Planck time becomes a composition depth: $\tP \leftrightarrow \nP \approx 56$ (for caesium, $d = 3$).
\end{enumerate}
\end{proposition}

\begin{proof}
(i) As derived in Proposition~\ref{prop:c_angular}. (ii) The energy $E = h\nu = \hbar \cdot 2\pi\nu$ distributed over $T(n,d)$ trajectories gives $\varepsilon = \hbar \cdot 2\pi\nu / T(n,d)$ per trajectory. (iii) By Definition~\ref{def:planck_depth} and Theorem~\ref{thm:caesium}, $\tP$ corresponds to the composition depth $\nP = 56$ at which the trajectory count first exceeds $\tau_{\mathrm{Cs}}/\tP$.
\end{proof}

\subsection{Implications for Measurement}

The composition-inflation mechanism reveals that temporal measurement resolution is not limited by clock speed but by the ability to distinguish categorical trajectories. These are fundamentally different resources:

\begin{proposition}[Resource Comparison]
\label{prop:resources}
\begin{enumerate}[label=(\roman*)]
\item \textbf{Clock speed} grows linearly with frequency: doubling $\nu$ doubles the tick rate. Resource cost: proportional to $\nu$ (energy, engineering, physical constraints).
\item \textbf{Trajectory distinguishability} grows exponentially with cycle count: each additional cycle multiplies the trajectory count by $d+1$. Resource cost: proportional to $n$ (observation time only).
\end{enumerate}
\end{proposition}

\begin{proof}
(i) Ticks per unit time $= \nu$; linear in $\nu$ by definition. (ii) $T(n+1,d) = (d+1) \cdot T(n,d)$ by Proposition~\ref{prop:growth}(ii); exponential in $n$ with base $d+1$. The cost of one additional cycle is one additional oscillator period $\tau = 1/\nu$, independent of the current trajectory count. Therefore, the trajectory count per unit of additional observation time grows exponentially, not linearly.
\end{proof}

This asymmetry between linear clock speed and exponential trajectory count is the physical essence of composition-inflation. Building a clock that ticks $10^{33}$ times per caesium period is impossible (it would require trans-Planckian clock speed). But observing 56 caesium ticks and counting the $3.8 \times 10^{33}$ categorical trajectories requires only $\sim 6$~ns of observation time and no exotic physics.


%==============================================================================
\section{Discussion}
\label{sec:discussion}
%==============================================================================

\subsection{Summary of Results}

We have derived the following:

\begin{enumerate}[label=(\arabic*)]
\item The composition-inflation formula: $T(n,d) = d \cdot (d+1)^{n-1}$ distinguishable categorical trajectories from $n$ oscillation cycles in $d$-dimensional entropy space (Theorem~\ref{thm:total_labeled}).

\item Angular resolution $\Delta\theta = 2\pi / T(n,d)$ is dimensionless and not subject to the Planck time constraint (Theorem~\ref{thm:no_planck}).

\item The Planck depth for caesium-133 is $\nP = 56$ cycles in $d = 3$ S-entropy space, corresponding to $3.8 \times 10^{33}$ distinguishable trajectories generated in 6.1~ns of physical time (Theorem~\ref{thm:caesium}).

\item All physically realizable oscillators reach Planck-equivalent angular resolution in 48--72 ticks, due to the logarithmic dependence of $\nP$ on oscillator frequency (Theorem~\ref{thm:universality}).

\item The five previously identified trans-Planckian enhancement mechanisms unify into composition-inflation at depth $n \approx 201$ (Theorem~\ref{thm:unification}).

\item Distance resolution via angular mapping exceeds the Planck length by 8 orders of magnitude for molecular-scale systems (Theorem~\ref{thm:distance_angle}).
\end{enumerate}

\subsection{What This Does Not Claim}

Precision about the scope of these results is essential:

\begin{enumerate}[label=(\roman*)]
\item We do \emph{not} claim to measure time intervals shorter than $\tP$. The composition-inflation mechanism produces categorical resolution, not temporal resolution. The distinction: ``how many seconds between events $A$ and $B$?'' has a Planck limit. ``How many distinguishable categorical trajectories separate events $A$ and $B$?'' does not.

\item We do \emph{not} claim to violate quantum gravity. The Planck scale remains the scale at which quantum gravitational effects dominate spacetime structure. Our claim is narrower: the Planck time constrains temporal intervals in SI units, not dimensionless categorical counts.

\item We do \emph{not} claim the Planck scale is wrong. We claim it is an artifact of the SI unit system's conflation of angular resolution (dimensionless) with temporal duration (seconds). In categorical angular units, the Planck time is simply a composition depth ($\nP \approx 56$), not a fundamental limit.

\item We do \emph{not} claim that any exotic physics is involved. The composition-inflation mechanism uses only: (a) integer compositions, a classical combinatorial object; (b) dimensional labeling in a finite-dimensional entropy space; and (c) the standard definition of angular resolution. All ingredients are well-established mathematics applied to bounded oscillatory systems.
\end{enumerate}

\subsection{Falsifiable Prediction}

If composition-inflation is physically realized, then a system with $n$ oscillation cycles in $d$-dimensional entropy space should exhibit $d \cdot (d+1)^{n-1}$ resolvable states, not $n$.

\begin{proposition}[Experimental Test]
\label{prop:test}
Consider a multi-mode oscillatory system with $d$ independent degrees of freedom, observed for $n$ oscillation cycles. The composition-inflation prediction is:
\begin{equation}
N_{\mathrm{resolved}} = d \cdot (d+1)^{n-1}.
\end{equation}
The linear prediction is $N_{\mathrm{resolved}} = n$. The ratio
\begin{equation}
R(n,d) = \frac{d \cdot (d+1)^{n-1}}{n}
\end{equation}
grows exponentially and exceeds $10^3$ for $n > 10$ ($d = 3$). This is a large, experimentally accessible signal.
\end{proposition}

\begin{proof}
$R(10,3) = 3 \cdot 4^9 / 10 = 3 \cdot 262{,}144 / 10 = 78{,}643 \approx 7.9 \times 10^4 > 10^3$. For $n > 10$, $R$ grows exponentially, making the distinction between exponential and linear trajectory counts measurable.
\end{proof}

The most direct experimental test would measure the number of distinguishable spectral features in a multi-mode oscillatory system (e.g., a trapped ion with multiple vibrational modes) as a function of observation time (cycle count). If the number of resolvable features grows exponentially with cycle count rather than linearly, composition-inflation is confirmed.

\subsection{Relation to Information Theory}

The composition-inflation mechanism has a natural information-theoretic interpretation. The information content of a trajectory specified by a labeled composition is:

\begin{proposition}[Information Content]
\label{prop:info}
The information required to specify one labeled composition from the set of $T(n,d)$ trajectories is
\begin{equation}
I(n,d) = \log_2 T(n,d) = \log_2 d + (n-1) \log_2(d+1) \;\;\text{bits}.
\end{equation}
For $d = 3$: $I(n,3) = \log_2 3 + 2(n-1) \approx 1.585 + 2(n-1)$ bits.
\end{proposition}

\begin{proof}
By Shannon's source coding theorem~\cite{shannon1948}, the minimum number of bits to specify one element from a set of $T$ equally likely elements is $\log_2 T$. Substituting $T = d \cdot (d+1)^{n-1}$ and using logarithm properties gives the formula.
\end{proof}

Each oscillation cycle contributes $\log_2(d+1)$ bits of information---2 bits per cycle for $d = 3$. This is more than the 1 bit per cycle that a simple binary counter would provide, reflecting the richer structure of labeled compositions compared to binary sequences.


%==============================================================================
\section{Angular Reformulation of Fundamental Constants}
\label{sec:angular_constants}
%==============================================================================

The composition-inflation mechanism shows that temporal resolution is angular, not durational. We now demonstrate that expressing fundamental constants in categorical angular units causes systematic cancellation, reducing physics to geometry plus a single irreducible coupling.

\subsection{Speed of Light as Angular Velocity}

The partition propagation bound~\cite{sachikonye2025ion} establishes that information propagates through bounded phase space at finite maximum speed:
\begin{equation}
c = \sup_n \, \ell_n \, \omega_n < \infty
\label{eq:c_partition}
\end{equation}
where $\ell_n$ is the spatial resolution at partition depth $n$ and $\omega_n$ is the maximum oscillation rate. This is derived from first principles: finiteness follows from finite mode density and finite oscillation rate at every level of the partition hierarchy~\cite{sachikonye2025ion}.

Now consider one complete oscillation cycle (one tick) of a bounded system with characteristic radius $r_0$. The oscillatory mode traverses one full phase cycle: $\Delta\phi = 2\pi$ radians. The distance traversed at the phase boundary is $\Delta x = 2\pi r_0$. By definition of one tick ($\Delta t = 1$ tick), the propagation speed is:
\begin{equation}
c = \frac{\Delta x}{\Delta t} = \frac{2\pi r_0}{1\;\text{tick}}
\label{eq:c_angular_distance}
\end{equation}

If we express distance in angular units ($\theta = x / r_0$), then:
\begin{equation}
\boxed{c = 2\pi \;\;\text{radians per tick}}
\label{eq:c_angular}
\end{equation}

This is not a unit convention but a theorem: the maximum partition propagation rate traverses exactly one full phase cycle per oscillation period, because one oscillation \emph{is} one phase cycle. The identity $c = 2\pi$ rad/tick is the angular expression of the partition propagation bound.

\subsection{Planck's Constant as Energy Per Radian}

The frequency-energy identity $E = \hbar\omega$ is derived from Bohr--Sommerfeld quantisation on bounded domains~\cite{sachikonye2025bounded}. For one tick ($\omega = 2\pi/\text{tick}$):
\begin{equation}
E_{\text{tick}} = \hbar \cdot \frac{2\pi}{\text{tick}} = \frac{h}{\text{tick}}
\label{eq:E_tick}
\end{equation}

Rearranging: $\hbar = E_{\text{tick}} / (2\pi)$. In categorical angular units:
\begin{equation}
\boxed{\hbar = E_{\text{tick}} / (2\pi) = \text{energy per radian}}
\label{eq:hbar_angular}
\end{equation}

This is exact: $\hbar$ is the energy carried by one radian of phase advance. A full tick carries $2\pi\hbar = h$ of energy.

\subsection{Mass as Angular Frequency Ratio}

From the ion trajectory completion framework~\cite{sachikonye2025ion}, inertial mass is:
\begin{equation}
m_0 = \frac{\hbar\omega_0}{c^2}
\label{eq:mass_from_ion}
\end{equation}
where $\omega_0$ is the localisation rest frequency. Substituting $\hbar = E_{\text{tick}}/(2\pi)$ and $c = 2\pi$ rad/tick:
\begin{equation}
m_0 = \frac{E_{\text{tick}} \cdot \omega_0}{2\pi \cdot (2\pi)^2} = \frac{E_{\text{tick}} \cdot \omega_0}{8\pi^3}
\label{eq:mass_angular}
\end{equation}

Since $\omega_0 = 2\pi\nu_0$ where $\nu_0$ is the rest frequency in ticks$^{-1}$:
\begin{equation}
\boxed{m_0 = \frac{E_{\text{tick}} \cdot \nu_0}{4\pi^2}}
\label{eq:mass_clean}
\end{equation}

Mass is the product of tick energy and rest frequency, divided by the geometric factor $(2\pi)^2$. The ``mysterious'' constants $\hbar$ and $c$ have dissolved into $2\pi$ factors and tick energy.

\subsection{The Mass-Energy Relation in Angular Units}

The rest energy $E_0 = m_0 c^2$~\cite{sachikonye2025ion} becomes:
\begin{equation}
E_0 = m_0 \cdot (2\pi)^2 = \frac{E_{\text{tick}} \cdot \nu_0}{4\pi^2} \cdot 4\pi^2 = E_{\text{tick}} \cdot \nu_0
\label{eq:E_mc2_angular}
\end{equation}

\begin{equation}
\boxed{E_0 = E_{\text{tick}} \cdot \nu_0 = \hbar\omega_0}
\label{eq:E0_clean}
\end{equation}

Rest energy equals tick energy times rest frequency. This is $E = \hbar\omega$ applied to the rest mode---the same formula as for any oscillation, confirming that $E = mc^2$ is not a separate postulate but the frequency-energy identity applied to the localisation mode.

\begin{figure*}[!htbp]
    \centering
    \includegraphics[width=\textwidth]{figures/panel_6_angular_constants.png}
    \caption{\textbf{Angular Reformulation of Fundamental Constants.}
    (\textbf{A})~Seven fundamental constants plotted in three-dimensional constant space with axes $\log_{10}(\text{SI value})$, categorical complexity (number of non-trivial factors in the angular expression), and constant index. Teal spheres indicate constants that are fully resolved in categorical angular units ($c$, $\hbar$, $k_B$, $m_0$, $E_0$, $t_P$); the red sphere marks $G$ as the sole irreducible constant. Sphere size scales with the number of $2\pi$ factors in the angular expression. The spatial separation of $G$ from the resolved cluster visualises its unique status as the only constant not reducible to geometry plus counting.
    (\textbf{B})~Constant reduction waterfall showing how each fundamental constant decomposes into categorical angular components. Horizontal stacked bars decompose each constant into geometric factors ($2\pi$, teal), tick energy ($E_{\text{tick}}$, blue), rest frequency ($\nu_0$, orange), and the gravitational constant ($G$, red). The speed of light $c = 2\pi$~rad/tick is purely geometric (single teal bar). Planck's constant $\hbar = E_{\text{tick}}/(2\pi)$ combines tick energy with one geometric factor. The Planck time $t_P = \sqrt{E_{\text{tick}} \cdot G}/(2\pi)^3$ requires three geometric factors plus the irreducible $G$ coupling. Bar lengths represent the $\log_{10}$ contribution of each component.
    (\textbf{C})~Count of $2\pi$ factors in the categorical angular expression of each constant. $c$ contains one factor (the phase cycle per tick). $\hbar$ contains one factor (the $2\pi$ relating $h$ to $\hbar$). $m_0$ contains two factors (from $c^2 = (2\pi)^2$). $t_P$ contains three factors (from $c^{5/2}$). $k_B$, $E_0$, and $G$ contain zero $2\pi$ factors---$k_B$ and $E_0$ because they involve $\ln(d+1)$ and $\nu_0$ rather than phase geometry, and $G$ because it is irreducible. The total geometric content of the Planck time (three factors of $2\pi$) reflects its construction from $c$ and $\hbar$, both of which carry $2\pi$.
    (\textbf{D})~First-principles derivation chain from the Bounded Phase Space Law to the irreducibility of $G$, displayed as a directed graph. Fourteen steps proceed through four categories: the axiom (gold), derived physics (teal: oscillatory necessity, partition coordinates, 3D space, speed of light $c$, Lorentz invariance, rest frequency $\omega_0$, mass $m_0$, $E = mc^2$), composition-inflation results (blue: $T(n,d)$, angular resolution), angular reformulation (purple: angular constants, Planck depth $n_P = 56$), and the terminal node (red: $G$ irreducible). Edges encode logical dependence. The chain requires no postulates beyond the axiom and no fundamental constants beyond $G$.}
    \label{fig:angular_constants}
\end{figure*}

\subsection{The Planck Time in Categorical Angular Units}

The Planck time is $t_P = \sqrt{\hbar G / c^5}$~\cite{planck1901}. Substituting $\hbar = E_{\text{tick}}/(2\pi)$ and $c = 2\pi$ rad/tick:
\begin{equation}
t_P = \sqrt{\frac{E_{\text{tick}} \cdot G}{2\pi \cdot (2\pi)^5}} = \frac{1}{(2\pi)^3}\sqrt{\frac{E_{\text{tick}} \cdot G}{2\pi}}
\label{eq:tP_angular}
\end{equation}

The factors of $2\pi$ are geometric (arising from the circular phase structure of oscillations). The irreducible content is:
\begin{equation}
\boxed{t_P^2 \propto E_{\text{tick}} \cdot G}
\label{eq:tP_irreducible}
\end{equation}

The Planck time is set by two quantities: (i) the energy per tick $E_{\text{tick}}$, which is the quantum of categorical action, and (ii) the gravitational constant $G$, which couples mass (= accumulated partition residue~\cite{sachikonye2025ion}) to spacetime curvature. Everything else is $2\pi$.

\begin{remark}
The gravitational constant $G$ is the only fundamental constant that does not reduce to a geometric factor or a tick-counting quantity in the categorical angular framework. This suggests that $G$ encodes the single irreducible physical content of the theory: the coupling between categorical structure (partition depth $\to$ mass) and spatial geometry (curvature). All other ``fundamental constants'' ($c$, $\hbar$, $k_B$) are artifacts of the SI unit system's use of seconds, meters, kilograms, and kelvins rather than ticks, radians, tick-energies, and categorical states.
\end{remark}

\subsection{Boltzmann's Constant as Entropy Per Trajectory}

Temperature in the bounded phase space framework is the categorical actualisation rate~\cite{sachikonye2025bounded}. The thermal energy per degree of freedom is $E_{\text{th}} = k_B T$. In categorical terms, $T$ is the tick rate (ticks per second), and $k_B T = k_B \cdot (\text{ticks}/\text{s})$. Using the composition-inflation count $T(n,d)$ for the number of distinguishable trajectories:
\begin{equation}
\boxed{k_B = \frac{E_{\text{tick}}}{\ln(d+1)}}
\label{eq:kB_angular}
\end{equation}

where $\ln(d+1)$ is the natural logarithm of the base of the composition-inflation exponential. For $d = 3$: $k_B = E_{\text{tick}} / \ln 4$. This identifies Boltzmann's constant as the energy cost per unit of categorical trajectory information---one nat of composition entropy.

\subsection{The Categorical SI Second}

With $c = 2\pi$ rad/tick, the SI second defined by caesium is:
\begin{equation}
1\;\text{second} = 9{,}192{,}631{,}770\;\text{ticks}
\label{eq:si_second}
\end{equation}

Each tick is one complete $2\pi$-radian phase cycle. The second is therefore:
\begin{equation}
1\;\text{second} = 9{,}192{,}631{,}770 \times 2\pi\;\text{radians of caesium phase}
\label{eq:second_radians}
\end{equation}

The SI meter follows from $c = 299{,}792{,}458$ m/s:
\begin{equation}
1\;\text{meter} = \frac{299{,}792{,}458}{9{,}192{,}631{,}770} \times 2\pi \approx 204.89\;\text{radians of caesium phase-distance}
\label{eq:meter_radians}
\end{equation}

The SI kilogram follows from $h = 6.626 \times 10^{-34}$ J$\cdot$s:
\begin{equation}
1\;\text{kg} = \frac{h \cdot \nu_{\text{Cs}}^2}{c^2} \times \left(\frac{1}{m_{\text{ref}}}\right)
\label{eq:kg_angular}
\end{equation}

In categorical units, every SI quantity reduces to (i) integer tick counts, (ii) factors of $2\pi$, (iii) the single irreducible constant $G$.

\subsection{Summary: Constant Reduction}

\begin{center}
\begin{tabular}{lcc}
\toprule
\textbf{Constant} & \textbf{SI expression} & \textbf{Categorical angular expression} \\
\midrule
$c$ & $299{,}792{,}458\;\text{m/s}$ & $2\pi\;\text{rad/tick}$ \\
$\hbar$ & $1.055 \times 10^{-34}\;\text{J}\cdot\text{s}$ & $E_{\text{tick}} / (2\pi)$ \\
$k_B$ & $1.381 \times 10^{-23}\;\text{J/K}$ & $E_{\text{tick}} / \ln(d+1)$ \\
$m_0$ & $\hbar\omega_0/c^2$ & $E_{\text{tick}} \cdot \nu_0 / (4\pi^2)$ \\
$E_0$ & $m_0 c^2$ & $E_{\text{tick}} \cdot \nu_0$ \\
$t_P$ & $\sqrt{\hbar G/c^5}$ & $\sqrt{E_{\text{tick}} G} / (2\pi)^3$ \\
$G$ & $6.674 \times 10^{-11}\;\text{m}^3\text{kg}^{-1}\text{s}^{-2}$ & \textbf{irreducible} \\
\bottomrule
\end{tabular}
\end{center}

Six of the seven fundamental constants reduce to tick counts, $2\pi$ factors, tick energy, and dimensionality $d$. Only $G$ remains irreducible---the sole input from gravity into the partition framework. This is the constant cancellation the categorical angular reformulation achieves: the Planck time, previously constructed from three independent constants ($\hbar$, $G$, $c$), depends on only one ($G$), with all other contributions being geometric ($2\pi$) or definitional ($E_{\text{tick}}$).

\begin{figure*}[!htbp]
    \centering
    \includegraphics[width=\textwidth]{figures/panel_7_categorical_si.png}
    \caption{\textbf{Categorical SI Units and the Irreducibility of $G$.}
    (\textbf{A})~Three-dimensional reduction graph showing the seven SI base units (second, metre, kilogram, ampere, kelvin, mole, candela) as faded nodes with edges converging toward five categorical primitives: ticks (integer count), $2\pi$ (geometric phase), $E_{\text{tick}}$ (energy quantum), $d$ (entropy space dimensionality), and $G$ (gravitational coupling). The convergence illustrates the unit reduction: all seven SI quantities are expressible in terms of these five categorical primitives. The gravitational constant $G$ stands apart as the only primitive that is not geometric or counting-based.
    (\textbf{B})~One SI second expressed as accumulated caesium phase. The $x$-axis shows tick count from $10^0$ to $10^{10}$ on logarithmic scale; the $y$-axis shows accumulated phase in units of $2\pi$ (turns). The teal line rises linearly: $\phi(n) = n$ turns $= 2\pi n$ radians. One second is reached at $n = 9{,}192{,}631{,}770$ ticks, corresponding to $5.775 \times 10^{10}$ radians of caesium phase. Reference lines at 1\,$\mu$s ($\sim$9200 ticks), 1\,ms ($\sim$9.2$\times 10^6$ ticks), and 1\,s mark temporal milestones. The linearity confirms that the second is a pure tick count---no additional structure is needed.
    (\textbf{C})~Decomposition of the Planck time $\log_{10}(t_P) = -43.27$ into contributions from each factor in the expression $t_P = \sqrt{\hbar G / c^5}$. The $\sqrt{\hbar}$ contribution ($-16.99$), $\sqrt{G}$ contribution ($-5.09$), and $c^{-5/2}$ contribution ($-21.19$) sum to the total. The categorical sub-decomposition further breaks $\sqrt{\hbar}$ into $\sqrt{E_{\text{tick}}}$ and $1/\sqrt{2\pi}$, and $c^{-5/2}$ into $(2\pi)^{-5/2}$ and frequency factors. The dominant contribution is from $c^{-5/2}$ (the speed-of-light suppression), which is purely geometric ($2\pi$ factors) in categorical units. The irreducible physics resides entirely in the $\sqrt{G}$ bar.
    (\textbf{D})~Irreducibility scatter: $\log_{10}(\text{SI value})$ versus $\log_{10}(\text{categorical residual})$ for each fundamental constant. The categorical residual is what remains after dividing out all factors of $2\pi$, $E_{\text{tick}}$, and $\nu_0$. Resolved constants ($c$, $\hbar$, $k_B$, $m_e$, $E_0$, $t_P$) cluster in the teal-shaded band near residual $\approx 1$ (i.e., they are fully explained by geometry and counting). The gravitational constant $G$ (red diamond) lies far outside this band, with a residual of $\sim 10^{-10}$, confirming its irreducibility. This is the central result of the angular reformulation: $G$ is the only fundamental constant that cannot be reduced to phase geometry and categorical counting.}
    \label{fig:categorical_si}
\end{figure*}


%==============================================================================
\section{First-Principles Derivation Chain}
\label{sec:derivation_chain}
%==============================================================================

We summarise the complete chain from axiom to Planck depth, referencing only results derived from bounded phase space geometry:

\begin{enumerate}
\item \textbf{Axiom}: Bounded Phase Space Law~\cite{sachikonye2025bounded}.
\item \textbf{Oscillatory necessity}: Poincar\'{e} recurrence $+$ bounded measure $\Rightarrow$ oscillatory dynamics is the unique valid mode (Theorem~\ref{thm:oscillatory}).
\item \textbf{Partition coordinates}: $(n, l, m, s)$ with $C(n) = 2n^2$~\cite{sachikonye2025bounded}.
\item \textbf{Three-dimensional space}: Angular partition geometry $\Rightarrow$ $\mathbb{R}^3$ with $SO(3)$~\cite{sachikonye2025ion}.
\item \textbf{Partition propagation bound}: $c = \sup_n \ell_n\omega_n < \infty$~\cite{sachikonye2025ion}.
\item \textbf{Lorentz invariance}: Oscillatory universality $+$ $SO(3)$ $\Rightarrow$ Lorentz group with invariant speed $c$~\cite{sachikonye2025ion}. No light postulate.
\item \textbf{Localisation rest frequency}: Confinement cost $\Rightarrow$ $\omega_0 > 0$~\cite{sachikonye2025ion}.
\item \textbf{Inertial mass}: $m_0 = \hbar\omega_0/c^2$~\cite{sachikonye2025ion}.
\item \textbf{Mass-energy}: $E_0 = m_0c^2 = \hbar\omega_0$ (theorem, not postulate)~\cite{sachikonye2025ion}.
\item \textbf{Composition-inflation}: $T(n,d) = d \cdot (d+1)^{n-1}$ (Theorem~\ref{thm:total_labeled}).
\item \textbf{Angular resolution}: $\Delta\theta = 2\pi/T$ is dimensionless and Planck-unconstrained (Theorem~\ref{thm:no_planck}).
\item \textbf{Angular constants}: $c = 2\pi$ rad/tick, $\hbar = E_{\text{tick}}/(2\pi)$ (\S\ref{sec:angular_constants}).
\item \textbf{Planck depth}: $n_P = 56$ for caesium (Theorem~\ref{thm:planck_depth}).
\item \textbf{Irreducible content}: Only $G$ remains---all other constants are geometric or definitional (\S\ref{sec:angular_constants}).
\end{enumerate}

Every step follows deductively from the Bounded Phase Space Law. No quantum mechanics is postulated (it is derived at step~2--3). No special relativity is postulated (it is derived at step~6). No constants are assumed (they are derived or shown to be geometric at step~12). The only external input is $G$, which couples partition depth to spatial curvature.


%==============================================================================
\section{Conclusion}
\label{sec:conclusion}
%==============================================================================

We have derived the composition-inflation mechanism from first principles: the combinatorics of integer compositions (Theorem~\ref{thm:comp_count}) combined with dimensional labeling in $d$-dimensional entropy space (Theorem~\ref{thm:total_labeled}) yields the central formula
\begin{equation}
\boxed{T(n,d) = d \cdot (d+1)^{n-1}}
\end{equation}
for the number of categorically distinguishable trajectories generated by $n$ oscillation cycles.

This formula is:
\begin{itemize}
\item \textbf{Exact}: derived from the binomial theorem with no approximations.
\item \textbf{Parameter-free}: $n$ is the cycle count and $d$ is the space dimension; both are given by the physical setup.
\item \textbf{Exponential}: each additional cycle multiplies $T$ by $d+1$.
\item \textbf{Planck-unconstrained}: the angular resolution $\Delta\theta = 2\pi/T$ is dimensionless and not subject to the Planck time bound.
\end{itemize}

The key quantitative result: for caesium-133 in 3-dimensional S-entropy space,
\begin{equation}
\boxed{\nP = 56 \;\;\text{cycles} \;\;\Longrightarrow\;\; T = 3 \cdot 4^{55} = 3.81 \times 10^{33} \;\;\text{trajectories}}
\end{equation}
exceeding the $2.02 \times 10^{33}$ Planck-time intervals in one caesium period. This is achieved in 6.1~ns of physical time, using no exotic physics---only the combinatorial richness of trajectory structure within 56 ordinary oscillation cycles.

The angular reformulation of fundamental constants (\S\ref{sec:angular_constants}) reveals why the cancellation occurs: in categorical angular units, $c = 2\pi$ rad/tick (exact), $\hbar = E_{\text{tick}}/(2\pi)$ (energy per radian), and $k_B = E_{\text{tick}}/\ln(d+1)$ (entropy per trajectory). Six of seven fundamental constants reduce to tick counts, factors of $2\pi$, tick energy, and entropy space dimensionality. Only the gravitational constant $G$ remains irreducible---the sole external input coupling partition depth to spatial curvature.

The complete derivation chain (\S\ref{sec:derivation_chain}) proceeds from the Bounded Phase Space Law through 14 deductive steps to the Planck depth result, without postulating quantum mechanics (derived at step~2--3), special relativity (derived at step~6), or any fundamental constants (derived or shown geometric at step~12). The Planck time is not a fundamental limit on measurement. It is an artifact of the SI unit system's conflation of angular resolution with temporal duration. Composition-inflation resolves this conflation and establishes that trans-Planckian categorical resolution is a theorem of combinatorics grounded in first-principles physics, not a violation of any physical law.


%==============================================================================
% BIBLIOGRAPHY
%==============================================================================

\begin{thebibliography}{99}

\bibitem{poincare1890} H. Poincar\'{e}, ``Sur le probl\`{e}me des trois corps et les \'{e}quations de la dynamique,'' \textit{Acta Mathematica} \textbf{13}, 1--270 (1890).

\bibitem{planck1901} M. Planck, ``\"{U}ber das Gesetz der Energieverteilung im Normalspectrum,'' \textit{Annalen der Physik} \textbf{309}(3), 553--563 (1901).

\bibitem{kac1947} M. Kac, ``On the notion of recurrence in discrete stochastic processes,'' \textit{Bulletin of the American Mathematical Society} \textbf{53}(10), 1002--1010 (1947).

\bibitem{shannon1948} C. E. Shannon, ``A mathematical theory of communication,'' \textit{Bell System Technical Journal} \textbf{27}(3), 379--423 (1948).

\bibitem{andrews1976} G. E. Andrews, \textit{The Theory of Partitions} (Addison-Wesley, Reading, MA, 1976).

\bibitem{heubach2009} S. Heubach and T. Mansour, \textit{Combinatorics of Compositions and Words} (CRC Press, Boca Raton, FL, 2009).

\bibitem{garay1995} L. J. Garay, ``Quantum gravity and minimum length,'' \textit{International Journal of Modern Physics A} \textbf{10}(2), 145--165 (1995).

\bibitem{bipm2019} Bureau International des Poids et Mesures, \textit{The International System of Units (SI)}, 9th edition (BIPM, S\`{e}vres, 2019).

\bibitem{sachikonye2025bounded} K. F. Sachikonye, ``On the thermodynamic consequences of bounded phase space: trajectory completion dynamics for geometric partitioning coordinates,'' Technical Report, Technical University of Munich (2025). See \texttt{https://github.com/fullscreen-triangle/lavoisier}.

\bibitem{sachikonye2025trans} K. F. Sachikonye, ``On the thermodynamic consequences of categorical state counting: oscillatory categorical partitioning based recursive harmonic networks,'' Technical Report, Technical University of Munich (2025). See \texttt{https://github.com/fullscreen-triangle/lavoisier}.

\bibitem{sachikonye2025ternary} K. F. Sachikonye, ``Categorical state counting in bounded phase space: digital mass spectrometry from partition dynamics,'' Technical Report, Technical University of Munich (2025). See \texttt{https://github.com/fullscreen-triangle/stella-lorraine}.

\bibitem{sachikonye2025ion} K. F. Sachikonye, ``On the geometric consequences of bounded phase space: partition depth trajectory dynamics in geometric coordinate systems,'' Technical Report, Technical University of Munich (2025). See \texttt{https://github.com/fullscreen-triangle/lavoisier}.

\end{thebibliography}

\end{document}